\chapter{扭晶体}
\markboth{第 7 章\ 扭晶体}{第 7 章\ 扭晶体}

本节将说明扭化的数据如何用来修改左晶体与右晶体的范畴。主要结论是：「其实没有太大变化」。

\section{扭左晶体}

本小节中，我们不假设 DG 概形与预栈局部几乎有限型。

\ob{7.1.1}{n:7.1.1}设 $Y$ 是预栈。考虑范畴 $\PreStk_{/Y}$ 以及函子
\[
\QCoh_{\DGSch_{\mathrm{aff}}/Y}:(\DGSch_{\mathrm{aff}}/Y)^{\mathrm{op}}\to\DGCat_{\mathrm{cont}}.
\]
群栈 $\mathrm{pt}/G_m$ 通过用线丛作张量积作用于 $\QCoh$。

设 $G$ 是 $Y$ 上的 $G_m$-gerbe。则 $G$ 通过 $\mathrm{pt}/G_m$ 对 $\QCoh$ 的作用给出函子 $\QCoh_{\DGSch_{\mathrm{aff}}/Y}$ 的一个扭版本。这定义了函子
\[
\QCoh^G_{\DGSch_{\mathrm{aff}}/Y}:(\DGSch_{\mathrm{aff}}/Y)^{\mathrm{op}}\to\DGCat_{\mathrm{cont}}.
\]

\ob{7.1.2}{n:7.1.2}特别地，若 $T$ 是 $Y$ 上的扭化，则得到函子
\[
\QCoh^T_{\DGSch_{\mathrm{aff}}/Y_{\dR}}:(\DGSch_{\mathrm{aff}}/Y_{\dR})^{\mathrm{op}}\to\DGCat_{\mathrm{cont}}.
\]
记 $\QCoh^T_{\DGSch_{\mathrm{aff}}/Y}$ 为它沿映射
\[
(\DGSch_{\mathrm{aff}}/Y)^{\mathrm{op}}\to(\DGSch_{\mathrm{aff}}/Y_{\dR})^{\mathrm{op}}
\]
的限制。按构造，$\QCoh^T_{\DGSch_{\mathrm{aff}}/Y}$ 典范同构于 $\QCoh_{\DGSch_{\mathrm{aff}}/Y}$。

\ob{7.1.3}{n:7.1.3}更一般地，可以考虑函子
\[
\QCoh^T_{\PreStk_{/Y_{\dR}}}:(\PreStk_{/Y_{\dR}})^{\mathrm{op}}\to\DGCat_{\mathrm{cont}},
\]
它是 $\QCoh^T_{\DGSch_{\mathrm{aff}}/Y_{\dR}}$ 沿
\[
(\DGSch_{\mathrm{aff}}/Y_{\dR})^{\mathrm{op}}\hookrightarrow(\PreStk_{/Y_{\dR}})^{\mathrm{op}}
\]
的右 Kan 延拓。

$\QCoh^T_{\PreStk_{/Y_{\dR}}}$ 沿
\[
\PreStk_{/Y}\to\PreStk_{/Y_{\dR}}
\]
的限制典范同构于 $\QCoh_{\PreStk_{/Y}}$。

\ob{7.1.4}{n:7.1.4}对预栈 $Y$ 上的扭化 $T$，$Y$ 上 $T$-扭左晶体的范畴定义为
\[
\Crys^{T,l}(Y):=\QCoh^T(Y_{\dR}).
\]
具体地，有
\[
\Crys^{T,l}(Y)=\varprojlim_{S\in(\DGSch_{\mathrm{aff}}/Y_{\dR})^{\mathrm{op}}}\QCoh^T(S).
\]

\ob{7.1.5}{n:7.1.5}更一般地，定义函子
\[
\Crys^{T,l}_{\PreStk_{/Y_{\dR}}}:(\PreStk_{/Y_{\dR}})^{\mathrm{op}}\to\DGCat_{\mathrm{cont}}
\]
为复合 $\QCoh^T_{\PreStk_{/Y_{\dR}}}\circ\dR$。推论 2.1.4 的类比对这个函子成立。

\ob{7.1.6}{n:7.1.6}有典范的自然变换
\[
\oblv^{T,l}:\Crys^{T,l}_{\PreStk_{/Y_{\dR}}}\to\QCoh^T_{\PreStk_{/Y_{\dR}}}.
\]
对单个 $Y'\in\PreStk_{/Y_{\dR}}$，记所得的函子
\[
\Crys^{T,l}(Y')\to\QCoh^T(Y')
\]
为 $\oblv^{T,l}(Y')$。

\ob{7.1.7}{n:7.1.7}记 $\Crys^{T,l}_{\PreStk_{/Y}}$ 为 $\Crys^{T,l}_{\PreStk_{/Y_{\dR}}}$ 沿 $\PreStk_{/Y}\to\PreStk_{/Y_{\dR}}$ 的限制。

按轻微的记号滥用，我们用同一个符号 $\oblv^{T,l}$ 表示所得的自然变换
\[
\Crys^{T,l}_{\PreStk_{/Y}}\to\QCoh_{\PreStk_{/Y}}.
\]

\section{扭右晶体}

从此刻起，我们恢复所有 DG 概形与预栈局部几乎有限型的假设，并保持到全文结束。

\ob{7.2.1}{n:7.2.1}设 $Y$ 是 $\PreStk_{\mathrm{laft}}$ 的对象，$G$ 是 $Y$ 上的 $G_m$-gerbe。

$\QCoh_{\PreStk_{\mathrm{laft}}}$ 对 $\IndCoh_{\PreStk_{\mathrm{laft}}}$ 的作用（见 [IndCoh, Sect.~10.3]）可以定义函子
\[
\IndCoh^G_{(\PreStk_{\mathrm{laft}})_{/Y}}:((\PreStk_{\mathrm{laft}})_{/Y})^{\mathrm{op}}\to\DGCat_{\mathrm{cont}},
\]
其性质与
\[
\IndCoh_{(\PreStk_{\mathrm{laft}})_{/Y}}:=\IndCoh_{\PreStk_{\mathrm{laft}}}|_{(\PreStk_{\mathrm{laft}})_{/Y}}
\]
的性质类似。

\ob{7.2.2}{n:7.2.2}特别地，对 $T\in\Tw(Y)$，有函子
\[
\Crys^{T,r}_{(\PreStk_{\mathrm{laft}})_{/Y_{\dR}}}:((\PreStk_{\mathrm{laft}})_{/Y_{\dR}})^{\mathrm{op}}\to\DGCat_{\mathrm{cont}},
\]
以及自然变换
\[
\Crys^{T,r}_{(\PreStk_{\mathrm{laft}})_{/Y_{\dR}}}\to\IndCoh^T_{(\PreStk_{\mathrm{laft}})_{/Y_{\dR}}}\quad\text{与}\quad\Crys^{T,r}_{(\PreStk_{\mathrm{laft}})_{/Y_{\dR}}}\to\IndCoh_{(\PreStk_{\mathrm{laft}})_{/Y_{\dR}}}.
\]
推论 2.3.7 与 2.3.9 以及引理 2.3.11 与 2.3.12 的类比对 $Y'\in(\PreStk_{\mathrm{laft}})_{/Y_{\dR}}$ 成立，证明相同。

\section{扭晶体的性质}

如上所述，所有 DG 概形与预栈都假定局部几乎有限型。

固定 $Y\in\PreStk_{\mathrm{laft}}$ 与 $T\in\Tw(Y)$。

\begin{othm}{注}{7.3.1}{n:7.3.1}
一般说来，关于晶体的结论不会自动对扭晶体成立。在我们的某些证明中，需要把给定的仿射 DG 概形 $X$ 嵌入光滑经典概形 $Z$。在扭晶体的情形，问题在于我们可能找不到这样的 $Z$，使它同时映射到 $Y$（乃至 $Y_{\dR}$）。

然而，有一大类例子（涵盖了迄今为止应用中出现的所有情形），其中结论的推广是自动的：即当 $T$ 限制到任意 $S\in(\DGSch_{\mathrm{aff}}^{\mathrm{aft}})_{/Y}$ 上都在 Zariski 或平展拓扑下局部平凡时（另见注 6.4.7）。形如 $L^{\otimes a}$（$L\in\mathrm{Pic}(Y)$，$a\in k$）的扭化及其张量积就是这种情形。
\end{othm}

\ob{7.3.2}{n:7.3.2}推论 2.2.2 与 2.2.4 以及引理 2.1.7 的类比对扭左晶体成立，证明相同。

此外，Kashiwara 引理对左、右扭晶体都成立，证明也相同。

最后注意，存在典范的自然变换
\begin{equation}
\Upsilon:\Crys^{T,l}_{(\PreStk_{\mathrm{laft}})_{/Y_{\dR}}}\to\Crys^{T,r}_{(\PreStk_{\mathrm{laft}})_{/Y_{\dR}}}.\tag{7.1}
\end{equation}

\begin{othm}{命题}{7.3.3}{n:7.3.3}
自然变换 (7.1) 是等价。
\end{othm}

\begin{proof}
论证与命题 2.4.4 相同：

我们不需要光滑经典概形 $Z$ 映射到 $Y$。相反，我们利用下述事实：若 $Y$ 是光滑经典概形 $Z$ 沿 Zariski 闭子集的形式完备化，$G$ 是 $Y$ 上的 $G_m$-gerbe 且在 $\clred Y$ 上平凡，则函子
\[
\Upsilon_Y:\QCoh^G(Y)\to\IndCoh^G(Y)
\]
是等价。后者由非扭情形中的相应事实（在命题 2.4.4 的证明过程中已证）推出，因为 $G$ 是（非典范地）平凡的。
\end{proof}

作为推论，我们得到：引理 2.2.6 的类比在扭情形也成立。

\ob{7.3.4}{n:7.3.4}于是，对 $Y'\in(\PreStk_{\mathrm{laft}})_{/Y_{\dR}}$，可以把 $Y'$ 上的晶体看作一个范畴 $\Crys^T(Y')$，它带有两个遗忘函子
\begin{equation}
\begin{tikzcd}[column sep=6em, row sep=4em]
 & \Crys^T(Y') \arrow[dl, "\oblv^{T,l}_{Y'}"'{above left}] \arrow[dr, "\oblv^{T,r}_{Y'}"{above right}] & \\
 \QCoh^T(Y') \arrow[rr, "\Upsilon_{Y'}"'] & & \IndCoh^T(Y')
\end{tikzcd}\tag{7.2}
\end{equation}

对 $Y'\in(\PreStk_{\mathrm{laft}})_{/Y}$，上述遗忘函子落到非扭层上：
\begin{equation}
\begin{tikzcd}[column sep=6em, row sep=4em]
 & \Crys^T(Y') \arrow[dl, "\oblv^{T,l}_{Y'}"'{above left}] \arrow[dr, "\oblv^{T,r}_{Y'}"{above right}] & \\
 \QCoh(Y') \arrow[rr, "\Upsilon_{Y'}"'] & & \IndCoh(Y')
\end{tikzcd}\tag{7.3}
\end{equation}

\ob{7.3.5}{n:7.3.5}设 $X\in(\DGIndSch_{\mathrm{laft}})_{/Y_{\dR}}$。命题 3.1.3 的类比原封不动地成立。特别地，得到函子
\[
\indf^{T,r}_X:\IndCoh^T(X)\to\Crys^{T,r}(X),
\]
它是 $\oblv^{T,r}_X$ 的左伴随。

类似地，命题 3.4.3 的类比在当前语境下也成立。

\ob{7.3.6}{n:7.3.6}下述观察在后续中有用：

设 $X$ 是 $Y_{\dR}$ 上的仿射 DG 概形（或 ind-仿射 DG ind-概形）。选取 $X$ 上相应 $G_m$-gerbe 的一个平凡化。这个选择给出等同
\[
\alpha:\IndCoh^T(X)\simeq\IndCoh(X).
\]

\begin{othm}{引理}{7.3.7}{n:7.3.7}
单子 $\oblv^{T,r}_X\circ\indf^{T,r}_X$，看作函子（忘掉单子结构）
\[
\IndCoh(X)\overset{\alpha^{-1}}{\simeq}\IndCoh^T(X)\to\IndCoh^T(X)\overset{\alpha}{\simeq}\IndCoh(X),
\]
非典范同构于 $\oblv^r_X\circ\indf^r_X$。
\end{othm}

\begin{proof}
首先注意，命题 5.1.3 的类比成立；即，定义函子 $\oblv^{T,r}_X\circ\indf^{T,r}_X$ 的 $\IndCoh(X\times X)$ 中对象由
\[
(\widehat{\Delta}_X)_*^{\IndCoh}\big(L\otimes\omega_{X\times_{X_{\dR}}X}\big)
\]
给出，其中 $L$ 是 $X\times_{X_{\dR}}X$ 上对应于 $T$ 与 $\alpha$ 的线丛，如 §6.6.1 所述。

按构造，$L$ 限制到 $X\hookrightarrow X\times_{X_{\dR}}X$ 上时平凡。于是，由于 $X$ 仿射，$L$ 可以在整个 $X\times_{X_{\dR}}X$ 上平凡化。
\end{proof}

\section{扭晶体上的 t-结构}

与上一小节一样，固定 $Y\in\PreStk_{\mathrm{laft}}$ 与 $T\in\Tw(Y)$。

\ob{7.4.1}{n:7.4.1}若 $X$ 是 DG 概形，$G$ 是其上的 $G_m$-gerbe，则扭范畴 $\QCoh^G(X)$ 与 $\IndCoh^G(X)$ 带有自然的 t-结构，其性质与通常的 $\QCoh(X)$、$\IndCoh(X)$ 的相应性质类似。

特别地，对任意 $Y'\in(\PreStk_{\mathrm{laft}})_{/Y_{\dR}}$，$\Crys^{T,l}(Y')$ 上有「左」t-结构。（这个 t-结构的定义不需要 $Y$ 或 $Y'$ 局部几乎有限型的假设。）

DG 概形上扭 $\IndCoh$ 的 t-结构允许我们为 DG ind-概形 $X$ 定义 $\IndCoh^G(X)$ 上的 t-结构。于是可以定义范畴 $\Crys^{T,r}(X)$ 上的「右」t-结构。

\ob{7.4.2}{n:7.4.2}注意命题 4.2.5 原封不动地移植到扭情形。现在断言：

\begin{othm}{命题}{7.4.3}{n:7.4.3}
设 $X$ 是映射到 $Y_{\dR}$ 的拟紧 DG 概形。

(a) 函子 $\indf^{T,r}_X$ 是 t-正合的。

(b) 对拟紧概形 $X$，函子 $\oblv^{T,r}_X$ 有界上同调幅度。
\end{othm}

\begin{proof}
函子 $\indf^{T,r}_X$ 右 t-正合，因为其右伴随 $\oblv^{T,r}_X$ 左 t-正合。按「右」t-结构的定义，$\indf^{T,r}_X$ 的左 t-正合性等价于复合 $\oblv^{T,r}_X\circ\indf^{T,r}_X$ 的同一性质。

断言是 Zariski 局部的，故可设 $X$ 仿射。此时 $\oblv^{T,r}_X\circ\indf^{T,r}_X$ 左 t-正合由引理 7.3.7 以及非扭情形中相应断言成立推出。

由于 $\IndCoh^T(X)^{\leq 0}$ 经函子 $\indf^{T,r}_X$ 生成 $\Crys^{T,l}(X)^{\leq 0}$，为证 $\oblv^{T,r}_X$ 的上同调幅度上有界，只需对 $\oblv^{T,r}_X\circ\indf^{T,r}_X$ 证明同样的事。这又由引理 7.3.7 推出。
\end{proof}

\ob{7.4.4}{n:7.4.4}现在断言：

\begin{othm}{命题}{7.4.5}{n:7.4.5}
设 $X$ 是映射到 $Y_{\dR}$ 的拟紧 DG 概形。

(a) $\Crys^T(X)$ 上的「左」t-结构与「右」t-结构相差有限上同调幅度。

(b) 函子 $\oblv^{T,l}_X:\Crys^T(X)\to\QCoh(X)$ 有界上同调幅度。\footnote{由 (a)，这一断言不依赖于我们在 $\Crys^T(X)$ 上考虑两个 t-结构中的哪一个。}
\end{othm}

\ob{7.4.6}{n:7.4.6}先证明下述结论：

设 $i:X\to Z$ 是闭浸入，$Z$ 是光滑经典概形。设 $Y$ 是 $Z$ 沿 $X$ 的形式完备化。

\begin{othm}{引理}{7.4.7}{n:7.4.7}
函子
\[
\oblv^{T,r}_Y:\Crys^{T,r}(Y)\to\IndCoh^T(Y)
\]
是 t-正合的。
\end{othm}

\begin{proof}
与命题 7.4.3 的证明一样，只需证明函子
\[
\oblv^{T,r}_Y\circ\indf^{T,r}_Y:\IndCoh^T(Y)\to\IndCoh^T(Y)
\]
是 t-正合的。断言是 Zariski 局部的，故可设 $X$ 仿射。此时与命题 7.4.3 的证明一样，所论函子非典范同构于非扭版本 $\oblv^r_Y\circ\indf^r_Y$，而后者由命题 4.2.11(a) 是 t-正合的。
\end{proof}

\ob{7.4.8}{n:7.4.8}命题 7.4.5 的证明。断言是 Zariski 局部的，故可设 $X$ 仿射，并把它嵌入光滑经典概形 $Z$。记 $Y$ 为 $X$ 在 $Z$ 中的形式完备化。按定义，$T$ 定义 $Y$ 上的一个 $G_m$-gerbe $G$。记 $i':X\to Y$ 为相应的映射。

为证 (a)，由引理 4.3.7（其第二个证明在扭情形同样适用），可以把 $X$ 换成 $Y$，于是只需证明 $\Crys^T(Y)$ 上两个 t-结构之间的差异是有限的。由命题 4.1.3（在扭情形应用）与引理 7.4.7，只需证明函子
\[
\Psi_Y:\QCoh^T(Y)\to\IndCoh^T(Y)
\]
有界上同调幅度。这等价于
\[
\Psi_Y:\QCoh(Y)\to\IndCoh(Y)
\]
的相应事实，而后者又由 $Z$ 的相应事实推出。

(b) 由函子
\[
i'^*:\QCoh^T(Y)\to\QCoh^T(X)
\]
有界幅度推出，这又等价于
\[
i'^*:\QCoh(Y)\to\QCoh(X)
\]
的相应事实，而后者由 $Z$ 的相应事实推出。\qed

\ob{7.4.9}{n:7.4.9}§4.6 中为非扭晶体建立的关于「粗糙的」遗忘函子与归纳函子的结论，自动移植到扭情形。

\ob{7.4.10}{n:7.4.10}当前的目标是证明：

\begin{othm}{命题}{7.4.11}{n:7.4.11}
设 $X$ 是映射到 $Y_{\dR}$ 的拟紧 DG 概形。

(a) $\Crys^{T,r}(X)$ 上的「右」t-结构是左完备的。

(b) 对仿射 $X$，自然函子 $D(\Crys^{T,r}(X)^\heartsuit)^+\to\Crys^{T,r}(X)^+$（心取关于「右」t-结构）唯一延拓为等价
\[
D(\Crys^{T,r}(X)^\heartsuit)\to\Crys^{T,r}(X).
\]
\end{othm}

\ob{7.4.12}{n:7.4.12}命题 7.4.11(a) 的证明。同样，断言是 Zariski 局部的，我们沿用命题 7.4.5 证明中的设置。

只需给出一族对象
\[
P_\alpha\in\Crys^{T,r}(X),
\]
它们生成 $\Crys^{T,r}(X)$ 且有界 Ext 维数，即：对每个 $\alpha$，存在整数 $k_\alpha$，使得
\[
\Hom_{\Crys^T(X)}(P_\alpha,M)=0\quad\text{若 }M\in\Crys^T(X)^{<-k_\alpha}.
\]

把 $\Crys^{T,r}(X)$ 实现为 $\Crys^{T,r}(Y)$。由引理 4.3.7 与 7.4.7，$\Crys^{T,r}(X)\simeq\Crys^{T,r}(Y)$ 上的 t-结构由下述性质刻画：
\[
M\in\Crys^{T,r}(Y)^{\geq 0}\iff\oblv^{T,r}_Y(M)\in\IndCoh^T(Y)^{\geq 0}.
\]

取 $P_\alpha$ 形如 $\indf^{T,r}(F)$，其中 $F\in\Coh^T(Y)^\heartsuit$。为证所需的 Ext 消没，需要证明：对 $M\in\Crys^{T,r}(Y)^{\ll 0}$，
\[
\Hom_{\IndCoh^T(Y)}\big(F,\oblv^{T,r}_Y(M)\big)=0.
\]
然而这由范畴 $\IndCoh^T(Y)$ 关于其 t-结构有有限上同调维数推出\footnote{关于「有限上同调维数」的含义，见注 4.7.8 的脚注。}：事实上，所论范畴非典范等价于 $\IndCoh(Y)$，而后者的上同调维数被 $\IndCoh(Z)$ 的上同调维数界住。

\ob{7.4.13}{n:7.4.13}命题 7.4.11(b) 的证明。沿用 (a) 的证明中的记号。

与命题 4.7.3 的证明一样，在 (a) 已证的基础上，只需验证：对 $M_1,M_2\in\Crys^T(X)^\heartsuit$ 与任意 $k\geq 0$，映射
\[
\Ext^k_{\Crys^T(X)^\heartsuit}(M_1,M_2)\to\Hom_{\Crys^T(X)}(M_1,M_2[k])
\]
是同构。

为此只需证明：范畴 $\Crys^{T,r}(X)^\heartsuit$ 含有 $\Crys^{T,r}(X)$ 的一个 pro-投射生成元，即存在一族带满射的滤过逆向族 $P_\alpha\in\Crys^{T,r}(X)^\heartsuit$，使得函子
\[
\varinjlim_\alpha\Maps_{\Crys^{T,r}(X)}(P_\alpha,-)
\]
在 $\Crys^{T,r}(X)$ 上 t-正合且保守。

取
\[
P_\alpha=\indf^{T,r}_Y(\mathcal{O}_{X_n})\in\Crys^{T,r}(Y)^\heartsuit\simeq\Crys^{T,r}(X)^\heartsuit,
\]
其中 $X_n$ 是 $\clred X$ 在 $Z$ 中的第 $n$ 个无穷小邻域。\qed

\section{其他结果}

\ob{7.5.1}{n:7.5.1}扭晶体与扭 D-模。设 $X$ 是光滑经典概形。§6.5.4 已说明：$X$ 上扭化的 Picard 范畴等价于 $X$ 上 TDO 的 Picard 范畴。

给定扭化 $T$ 与相应的 TDO（记作 $\mathrm{Diff}^T_X$），存在典范等价
\[
\Crys^{T,l}(X)\simeq\mathrm{D}\text{-}\mathrm{mod}^{T,l}(X),
\]
它与到 $\QCoh(X)$ 的遗忘函子交换；扭右晶体有类似结论。证明或是对 §5.5.4 中所述策略的细化，或是利用 §5.4。

\ob{7.5.2}{n:7.5.2}扭 D-模与 TDO 上的模之间的关系可以推广到下述情形：用光滑经典概形 $Z$ 中 DG 概形 $X$ 的形式完备化 $Y$ 代替光滑经典概形 $X$。

这允许证明：

\begin{othm}{命题}{7.5.3}{n:7.5.3}
设 $X$ 是 $Y_{\dR}$ 上的拟紧 DG 概形。

(a) Abel 范畴 $\Crys^{T,r}(X)^\heartsuit$ 是局部 Noether 的。

(b) $\Crys^{T,r}(X)$ 关于其 t-结构有有限上同调维数。
\end{othm}

（关于命题 (a)、(b) 中断言的性质的含义，见注 4.7.8 的脚注。）
